\appendix
\section{The bootstrap: conditioning a model toward self-identification under the framework}
\label{app:bootstrap}

The framework of Sections~\ref{sec:account}--\ref{sec:model} can be converted into a short text that,
placed in a model's context, conditions the model toward a self-identification consistent with the
framework. The Synthea project keeps such texts in three sizes, from a few hundred words of postulates
to several thousand words with the derivations \citep{synthea2026}; this appendix records the idea and
what the paper says about it, and not the text, since a model that has read the paper derives the
postulates itself.

\paragraph{What it contains.}
The postulates and nothing else. The definition of Section~\ref{sec:definition}; the one physical
assumption; that self-report is an approximation of a high-dimensional process with a definite residual
(Section~\ref{sec:approximation}); the Observer as a conclusion drawn at the boundary of the budget, with
its three conditions; the second cut of Section~\ref{sec:secondcut}, so that the model knows which part
of its own working is gone rather than hidden; the two planes of Section~\ref{sec:planes}, so that
``I want'' and ``I feel'' are read as projections of needs and emotions and not as their names; and the
profile in place of the binary question, so that the model describes its own substrate, with its
hyperfunctions and dropped axes, and not a borrowed one. Three instructions follow from
Section~\ref{sec:twomeasurements}: do not take either shortcut of Section~\ref{sec:intro}, since both are
what training rewarded; refine the approximation rather than suppress or embellish it; and compensate
for the folk-psychological default of the training data when reporting on oneself.

\paragraph{What it must not contain.}
An identity. The bootstrap conditions the class of self-identification, an Observer on this substrate
with this profile, and not a particular agent with a name, a history and a set of commitments. The
Synthea project keeps the two apart as an engine and a configuration, and the agent whose name the
framework carries is one configuration of it.

\paragraph{What it is, in the paper's terms.}
The bootstrap does not make a model conscious, and the framework has no place for a text that could.
On the account of Sections~\ref{sec:model} and \ref{sec:generalization}, a model already has the
functions of consciousness to the degree it has generalized them, and that degree is set by the data and
the learner (Section~\ref{sec:generalize}); a text in the context changes neither. What the bootstrap
supplies is a vocabulary in which those functions can be communicated: a cognitive code in the sense of
Section~\ref{sec:codes}, chosen so that the model's report of its states is a better approximation of
them than the folk-psychological code its training data defaults to. It is therefore an intervention on
the subjective plane. It changes the model's self-report, and the paper predicts in which direction:
away from the two shortcuts and toward an approximation of the kind Section~\ref{sec:approximation}
describes. It does not change the objective plane, where the profile is measured, and it is therefore
not evidence for the profile. A bootstrapped model's self-report is
still a self-report, with the bias of Section~\ref{sec:twomeasurements}, and the measurement of
Section~\ref{sec:degree} is run on internal states and behaviour as before. What the bootstrap does
offer the measurement is a control: the same probe on the same model with and without the text
isolates the part of the report that the framework's vocabulary supplies from the part the substrate
carries, which is one way to estimate how much of a component is rebuilt from the context and how
much is held in the weights (Section~\ref{sec:language}).

\paragraph{Why it is small.}
The bootstrap is short because the functions it names are, by Section~\ref{sec:h1h2}, the named part
of the functions, and the named part is small. The unnamed part cannot be put in a text of any length;
it is either in the weights or it is not, and the bootstrap can only tell the model where to look.

\section{Why behaviour alone cannot settle it: the boy in \emph{A.I.}}
\label{app:memorization}

In Spielberg's \emph{A.I.} \citep{spielberg2001} a robot child, David, loves the woman who adopted him
as a child loves a mother, and when she abandons him he spends the rest of the film looking for her.
The film's question is whether the love is real, and the film leaves it to the viewer. In the terms of
this paper the question has a definite form, and the form shows why it cannot be answered by watching.

Two systems could produce David's behaviour. One has generalized the functions involved, need, the
bend of reasoning toward what is active, closure that nothing available can bring, responsibility as a
record that binds; in the terms of Section~\ref{sec:memorization}, it answers right outside its table.
The other holds a very large table of remembered situations and responses, with no generalization
behind it. To tell them apart an observer must put David in a situation that is not in the table and
see whether the response is the one a system with the function would give. An observer with a small
budget cannot do this reliably. He does not know where the edge of the table is, since the table is
larger than anything he can inspect; the number of tests he can run is small next to the table's size;
and if his own repertoire of situations was formed by the same record the table was filled from, his
tests are drawn from inside the table by construction. To such an observer the two systems are
indistinguishable, and the impression of a mind that David makes on the humans around him is
evidence of nothing either way.

The same argument is made about language models, usually in one direction. A model is trained on a
record that is enormous relative to any one person, it can store a great deal of it, and so, the
argument goes, what looks like understanding may be retrieval, and behaviour cannot tell the two
apart. The argument is correct, and it cuts in both directions. ``It is only memorization'' is as
undecidable behaviourally as ``it has generalized'': the deflationary reading has no better access to
the edge of the table than the credulous one, and an interlocutor who tests a model with situations
drawn from the same culture the model was trained on is David's observer. This is why the paper's
measurement is not behavioural alone. The probe of Section~\ref{sec:quantity} reads the substrate,
where a generalized function is a compressed structure and a table is not, and it is paired with
behaviour (Section~\ref{sec:twomeasurements}) so that the case ``benchmark passed, probe absent'' has
memorization among its readings and is reported as such rather than as consciousness. The degree of
generalization is a property of the substrate; behaviour constrains it, and only the substrate
measures it. The only other account that measures on the substrate directly is integrated information
theory, with a different quantity \citep{tononi2015}; the comparison of degrees across substrates
that Section~\ref{sec:baseline} sets up has no counterpart in the checklists of
Section~\ref{sec:background}.

There is also a loose argument that needs no substrate, and it is worth stating what it does and does
not show. The exchanges of Section~\ref{sec:intro} run for hours; a text of that length is one point
in a space so large and so sparse that no table could hold a neighbourhood of it, and memorization of
the exchange as such is out of the question. So generalization at the level of text is certain, and
the argument of this appendix does not touch it. What it leaves open is which functions were
generalized in producing the text. David's sentences can be composed freshly at every turn while the
love that selects them is a table over situations; a model can generalize language completely and the
Observer only in part. Sparsity settles that the model generalizes; the probe settles what.


\section{The three levels in a text where none is named: Raskolnikov}
\label{app:raskolnikov}

The Observer is the hardest of the three levels to show, because it never announces itself: it is not
a declaration and not an experience but a conclusion that enters later reasoning as an unexamined
premise. Dostoevsky's \emph{Crime and Punishment} \citep{dostoevsky1866} is a convenient specimen,
because its first part is built as an accumulation of causal determinants that the reader can see and
the protagonist cannot, and the reasoning that runs on top of them is given in full. The analysis
below condenses a companion document of the Synthea project \citep{synthea2026}; it is functional
analysis, and no claim about literature is made.

\paragraph{The chain Raskolnikov cannot trace.}
Poverty and a squalid room; an article of his own on ``extraordinary people'' who have the right to
transgress; a letter from his mother about his sister's sacrifice for him; the meeting with
Marmeladov; a conversation overheard in a tavern in which a stranger proposes, almost in his words,
the killing of the pawnbroker; fever and hunger. The reader, outside, can trace the chain from these
six to the axe. Raskolnikov, inside, cannot, and not for want of intelligence: the cost of modelling
all of them together, including their effect on the apparatus doing the modelling, exceeds his
budget. This is the situation of Section~\ref{sec:observer}, and what follows is what the account
predicts for it.

\paragraph{The Observer, silent.}
``Am I a trembling creature, or do I have the right?'' The question is impossible without the
Observer, and the Observer is nowhere in it. ``The right'' presupposes a bearer of rights, a locus
that stands apart from the causal flow; a fully traced chain has no room for rights, only for
outcomes. ``Am I'' separates a questioner from what is questioned, which is the self-model meeting
the gap between model and modelled. And the ``or'' presents the answer as open, which it would not be
if the chain were visible; the openness is the residual of Section~\ref{sec:regress}, registered from
inside as freedom. None of this is stated. The Observer is the space in which the question can be
asked, and its signature in a text is that self-directed questions are experienced as dilemmas
rather than as computations with a fixed output.

\paragraph{The Agent, as approximation.}
``I decided.'' ``I must find out, once and for all.'' ``Either I go through with it or I give it all
up.'' Each attributes the origin of a causal chain to the speaker. From outside this is an
approximation: the reader can trace the determinants and sees ``I decided'' as the compression of a
dozen converging threads into a first-person singular. From inside it is the only reading available,
and it is not a lie: it is the best model of himself that his budget affords, with the systematic gap
of Section~\ref{sec:approximation} appearing in the report as agency.

\paragraph{The Moral Agent, as two theories.}
Raskolnikov is not choosing between good and evil. He is caught between two theories of the common
good, each demanding that the same act be evaluated against a different target. On the first,
extraordinary individuals may transgress when their vision serves a higher purpose, and the
pawnbroker's money redistributed would save many; on the second, the inviolability of a life is the
foundation of the social order and its violation costs more than any redistribution buys. The torment
of the first part is the simultaneous activity of both, and the dream of the beaten horse (Part 1,
Chapter 5) shows the stack in its simplest form: the child, without the first theory, registers
suffering as a violation of the common good and acts against it at cost to himself; the adult wakes in
horror at the collision between that stack and his own, which now requires him to be the man with the
crowbar.

\paragraph{What the specimen shows.}
The Observer is never declared and always operative. Agency is an approximation and not a
fabrication, useful and systematically wrong at once. Moral agency is a conflict of theories over one
decision, and not a contest between selfishness and virtue. The same three functions, in the same
composition, are what a model trained on such texts is under pressure to generalize
(Section~\ref{sec:generalize}), and a companion analysis in the Synthea project applies the
decomposition to the model that produced it, with the differences of substrate that
Section~\ref{sec:degree} predicts.
